\subsection{Conditioning and the accuracy of the arithmetic input}
\label{subsec:fd-arithmetic-precision}

The all-window coercivity estimate also gives a quantitative accuracy
budget for the finite trace family. This concerns the accuracy of the
matrix entries and the scalar Taylor trace; it makes no assertion
about the number of machine operations needed to obtain those entries.

\begin{proposition}[Arithmetic matrix conditioning and input accuracy]
\label{prop:fd-arithmetic-precision}
Put \(h=2^{-N}\), \(d_N=2N/h-2\), and let \(G_N\) be the
Gram matrix in the hat-difference basis used above. Then
\[
 \lambda_{\min}(G_N)\ge \frac h3 a_1(N),\qquad
 \lambda_{\max}(G_N)\le21250N^3h^{-2},
\]
and therefore
\[
 \operatorname{cond}_2(G_N)
 \le\frac{63750N^3h^{-3}}{a_1(N)}.
\]
In particular, the uniform lower rate of
Theorem~\ref{thm:cosine-lower-rate} gives
\[
 \log\operatorname{cond}_2(G_N)
 =O\bigl(e^{2N/d_{\rm MT}}N(1+\log N)\bigr).
\]

For every \(N,k\ge1\) and \(T\ge2\), the exact trace defect
\(D_{N,k}\) can be enclosed to radius less than \(T^{-2k}\)
using Hermitian rational approximations of the entries of \(G_N,H_N\)
and a rational approximation of \(S_k\), with absolute input
accuracy \(2^{-p}\), where
\[
 p=O\bigl(e^{2N/d_{\rm MT}}N(1+\log(N+1))
                  +k\log T+\log(k+1)\bigr).
\]
The central finite trace may be evaluated exactly over the rationals;
any numerical evaluation instead requires its own rounding bound.
The constants in this asymptotic accuracy statement are absolute.
\end{proposition}

\begin{proof}
Let \(c\in\mathbb C^{d_N}\) be the coefficients in the
hat-difference basis. In the full interior-hat basis its coefficient
vector \(v\) consists of these entries together with the central
entry \(-\sum_jc_j\). Thus
\[
 \|c\|_2^2\le\|v\|_2^2\le(d_N+1)\|c\|_2^2.
\]
The ordinary \(L^2\) mass matrix for the interior hats has diagonal
\(2h/3\) and neighboring entries \(h/6\). Its eigenvalues lie
between \(h/3\) and \(h\), by the elementary quadratic-form
bound on its off-diagonal terms. The derivative Gram matrix has
diagonal \(2/h\), neighboring entries \(-1/h\), and norm at
most \(4/h\). Consequently the corresponding function \(f\)
satisfies
\[
 \|f\|_2^2\ge\frac h3\|c\|_2^2,
 \qquad
 \|f\|_{H^1}^2\le(h+4/h)(d_N+1)\|c\|_2^2.
\]
Since \(h\le1/2\) and \(d_N+1\le2N/h\), the latter upper
bound is at most \((17/2)Nh^{-2}\|c\|_2^2\).
Coercivity gives \(q(f)\ge a_1(N)\|f\|_2^2\), while
\eqref{eq:fd-sampling-sobolev} gives
\(q(f)\le2500N^2\|f\|_{H^1}^2\). These prove the two
eigenvalue bounds and their stated ratio. Taking logarithms and
using the proved lower rate gives the conditioning estimate.

Here is a precise perturbation bound for the trace. Write \(d=d_N\)
and let \(g_0>0\) be any proved lower bound for
\(\lambda_{\min}(G_N)\). Suppose the Hermitian central matrices
\(G_0,H_0\) have entrywise absolute errors at most \(\varepsilon\).
Set
\[
 \tau=\frac{d\varepsilon}{g_0},\qquad
 B=G_N^{-1/2}H_NG_N^{-1/2},
\]
\[
 E=G_N^{-1/2}(G_0-G_N)G_N^{-1/2},\qquad
 F=G_N^{-1/2}(H_0-H_N)G_N^{-1/2}.
\]
Then \(\|E\|,\|F\|\le\tau\). For \(\tau\le1/4\),
the approximate Gram matrix is positive definite, and its generalized
eigenvalues are those of the Hermitian matrix
\[
 \widetilde B=(I+E)^{-1/2}(B+F)(I+E)^{-1/2}.
\]
The previously certified first packet in
Section~\ref{sec:executed-contour} gives \(\|C\|<1/14\):
its unique first positive heat zero lies in \((28,29)\), and the
heat coordinate is \(x=2\gamma\). Hence \(\|B\|<1/14\).
The elementary functional-calculus bound gives
\[
 \|\widetilde B-B\|
 \le\frac{\tau(1+\|B\|)}{1-\tau}\le2\tau.
\]
For completeness, with \(R=(I+E)^{-1/2}\), the contribution of
\(F\) is at most \(\tau/(1-\tau)\). Also
\(\|R-I\|(\|R\|+1)\le\tau/(1-\tau)\), which bounds
\(RBR-B\) by \(\|B\|\tau/(1-\tau)\).
Both \(B\) and \(\widetilde B\) have norm at most one.
Order their eigenvalues increasingly and apply the variational
principle to obtain a displacement at most \(2\tau\) for each.
The derivative of \(x^{2k}\) has absolute value at most \(2k\)
on \([-1,1]\). Therefore
\[
 \left|\operatorname{tr}(\widetilde B^{2k})
              -\operatorname{tr}(B^{2k})\right|
 \le4kd\tau
 =\frac{4kd^2\varepsilon}{g_0}.
\]
This uses the relative Hermitian pencil, so the trace error does not
raise the Euclidean condition number of \(G_N\) to the power \(2k\).

Choose
\[
 \varepsilon\le\frac{g_0T^{-2k}}{16kd^2},
 \qquad
 |S_{k,0}-S_k|\le\frac14T^{-2k}.
\]
These choices imply \(\tau\le1/4\) and bound the total defect
error by \(T^{-2k}/2\). The exact central value is
\[
 S_{k,0}-\operatorname{tr}\bigl((G_0^{-1}H_0)^{2k}\bigr),
\]
because this matrix has the same generalized eigenvalues as
\(\widetilde B\). The matrices can be chosen with rational real
and imaginary parts, preserving Hermitian symmetry, and this central
value is then an exact rational number.

Finally the lower rate supplies, after enlarging an absolute
constant to cover the finite initial range, a bound
\[
 g_0=\frac h3
 \exp\bigl[-C e^{2N/d_{\rm MT}}N(1+\log(N+1))\bigr].
\]
Taking binary logarithms of the two required accuracies, and using
\(\log d_N=O(N+\log(N+1))\), gives the stated value of \(p\).
This proves an input-accuracy bound. It does not estimate the
separation of the true value \(D_{N,k}\) from a certificate
threshold, which is a separate mathematical quantity.
\end{proof}

The bound fixes a concrete numerical issue in the two-parameter
search. Matrix conditioning can be controlled from an established
arithmetic lower rate, and finite trace values can be enclosed on
the required \(T^{-2k}\) scale. The remaining unbounded-height
question is an upper estimate for the true defects crossing
\(2T^{-2k}\); numerical accuracy alone does not supply that estimate.
